\begin{table}[htb]
\begin{tabular}{p{8cm}p{8cm}}
{\Large \bfseries Resumen} &{\Large \bfseries Abstract} \\ \\

Bajo el marco de trabajo PRISMA, se realizó una revisión sistemática de la literatura en la base de datos Scopus, donde se analizaron 27 artículos publicados entre los años 2022 y 2026. Los resultados revelaron que si bien modelos multilingües no se desempeñan correctamente en lenguas de bajos recursos, por ejemplo, el modelo \textit{Cross-Lingual Language Modeling-Roberta}(XLM-R) en un escenario \textit{zero-shot} sobre 10 lenguas indígenas obtuvo una presición promedio de 38.48\%, existen enfoques que permiten obtener resultados más favorables, como el uso de meta-aprendizaje (MAML), implementación de arquitecturas hibridas, inyección de conocimiento lingüistico y un enfoque participativo con la población de estudio. Aún existen limitaciones que son necesarias de abarcar en futuras investigaciones, como la escacez de datos y falta de recursos computacionales, la necesidad de un enfoque participativo y métricas especializadas para lenguas de bajos recursos lingüisticos, por tal motivo se sugiere que los investigadores adopten enfoques participativos y prioricen las necesidades con los hablantes de las lenguas estudiadas.

&Under the PRISMA methodology, a systematic literature review was conducted using the Scopus database, in which 27 articles published between 2022 and 2026 were analyzed. The results revealed that, although multilingual models do not perform adequately in low-resource languages—for example, the \textit{Cross-Lingual Language Modeling-Roberta} (XLM-R) model achieved an average accuracy of 38.48\% in a \textit{zero-shot} scenario across 10 Indigenous languages—there are approaches that make it possible to obtain more favorable results, such as the use of meta-learning (MAML), the implementation of hybrid architectures, the injection of linguistic knowledge, and a participatory approach involving the study population. There are still limitations that need to be addressed in future research, such as data scarcity and the lack of computational resources, the need for a participatory approach, and specialized metrics for low-resource languages. Therefore, it is suggested that researchers adopt participatory approaches and prioritize needs together with the speakers of the languages being studied.
% Debe contener un máximo de 230 palabras, no debe incluir ecuaciones o referencias. El contenido del resumen debe estar completamente justificado.
……………..
\\[-0.3cm]
\palabrasclave{Aprendizaje, aprendizaje personalizado, enseñanza, equidad educativa, Inteligencia Artificial, lenguas con pocos recursos, pedagogía, procesamiento del lenguaje natural.}&\keywords{Artificial Intelligence, educational equity, learning, low-resource languages, natural language processing, pedagogy, personalized learning, teaching.}  \\
\end{tabular}
\end{table}

\newpage
